\section{A filtered degeneration of the pentagon}
\label{sec:degeneration}

We now work over $\F$.  The elements to which we apply the
argument are the reductions $\bar\psi\in\overline L_n$ of
Proposition~\ref{prop:saturated-reduction}, with $n>2$.  In particular,
they are ordinary Lie polynomials and satisfy the reduced pentagon.
We first introduce a filtration in which a leading part of that pentagon
can be tested in a different Lie algebra.

\begin{samepage}
\begin{lemma}\label{lem:abc-free}
The assignments
\[
 A\longmapsto x^2,\qquad C\longmapsto[x,y],\qquad B\longmapsto y
\]
define an injective homomorphism
$\F\langle A,C,B\rangle\longrightarrow\F\langle x,y\rangle$.
Its restriction embeds the ordinary free Lie algebra
$\Lie_{\F}\langle A,C,B\rangle$ in the latter associative algebra.
Every homogeneous element of $\Lie_{\F}\langle x,y\rangle$ of weight
greater than one belongs to this embedded Lie algebra.
\end{lemma}
\end{samepage}

\begin{proof}
Order words of each fixed length lexicographically with $x>y$.
The leading words of $x^2,[x,y],y$ are respectively $xx,xy,y$, each with
coefficient one.  These form a prefix code: none is a proper initial
segment of another.  Thus their concatenations have unique decodings,
so distinct words in $A,C,B$ have distinct leading words.  Split a
putative relation by its total $x,y$ length and take the largest
leading word in each part; its coefficient proves associative
injectivity.  Ordinary free Lie algebras embed in their free associative
enveloping algebras, so the Lie assertion follows as well.

We use the ordinary free-Lie elimination construction
\cite[Theorem~0.6]{Reutenauer}, whose needed instance can be seen directly.
Set $e_j=\ad_x^j y$ for $j\geq0$.  The Lie subalgebra generated by the
$e_j$ is stable under $\ad_x$, because $\ad_x(e_j)=e_{j+1}$ and
$\ad_x$ is a derivation.  Its sum with the line $\F x$ is therefore a
Lie algebra containing $x,y$, hence the whole free Lie algebra.  The
subalgebra is thus an ideal with one-dimensional quotient spanned by
the class of $x$.  Every homogeneous Lie polynomial of weight greater
than one consequently lies in this subalgebra.  It is itself free on
the $e_j$: their leading words $x^j y$ form a prefix code, which proves
injectivity of the corresponding free associative algebra and hence
of its ordinary free Lie algebra.

In characteristic two, the associative identity
$\ad_x^2(P)=[x^2,P]$ gives
\begin{equation}\label{eq:elimination-abc}
 e_{2i}=\ad_A^i B,\qquad e_{2i+1}=\ad_A^i C\qquad(i\geq0).
\end{equation}
Both expressions belong to the ordinary Lie algebra on the three
abstract letters $A,C,B$, which proves the final assertion.
\end{proof}

The statement does not assert that $x^2$ is an ordinary Lie polynomial
in $x,y$.  Rather, it embeds an ordinary Lie algebra on three new
letters into an associative algebra; the original Lie polynomials of
weight greater than one happen to lie in its image.

Give $A,C$ weight two and $B$ weight one, and let $F^r$ be the span of
ordinary Lie monomials containing at least $r$ occurrences of $B$.
The multigrading of the free Lie algebra makes this a decreasing
filtration.  Through Lemma~\ref{lem:abc-free}, we regard
$\overline L_n$ as a subspace of
$\Lie_{\F}\langle A,C,B\rangle$ and give it the induced filtration
$F^r\overline L_n=\overline L_n\cap F^r$; write
$\operatorname{gr}^r_F\overline L_n=
F^r\overline L_n/F^{r+1}\overline L_n$.
Write $\Psi(A,C,B)$ for the unique polynomial representing
$\bar\psi$ under Lemma~\ref{lem:abc-free}.  If $\Psi\in F^r$, let
$\chi(A,C,B)$ be its component containing exactly $r$ letters $B$.
Whenever $\chi\ne0$, there is an integer $m\geq0$ such that
\begin{equation}\label{eq:leading-count-weight}
 n=2m+r,
\end{equation}
and every associative word of $\chi$ has $m$ letters from $\{A,C\}$
and total length $m+r$.  The case $m=0$ is zero: an ordinary Lie
polynomial in $B$ alone has only its degree-one component, whereas
$n>2$.  Thus a nonzero leading part always has $m>0$, and in particular
$n>r$.  We shall later study the ordinary Lie polynomial
$f=\chi(0,C,B)$.  This section establishes the relation satisfied by
$\chi$ before that projection.

We need two algebras with the same named symbols.  The source is the
ordinary enveloping algebra of the infinitesimal braid Lie algebra
on strands $0,1,\ldots,d$, over $\F$.  In it put
\begin{equation}\label{eq:source-symbols}
\begin{gathered}
 X_i=t_{0i},\qquad u_{ij}=t_{ij},\qquad
 a_i=X_i^2,\\
 v_{ij}=[X_i,u_{ij}]=[X_j,u_{ij}]=[X_i,X_j].
\end{gathered}
\end{equation}
The equalities defining $v_{ij}$ follow from the three-strand braid
relations and characteristic two.  Here the $a_i$ are associative
squares.

For the target, let $\mathcal A_d$ be the arrangement in $\Q^d$ with
hyperplanes
\[
 \begin{gathered}
 z_i=0\quad(1\leq i\leq d),\\
 z_i-z_j=0,\qquad z_i+z_j=0\quad(1\leq i<j\leq d).
 \end{gathered}
\]
Let $\mathfrak h_{\Z,d}$ be the ordinary Lie algebra over $\Z$ with
one generator $t_H$ for each $H\in\mathcal A_d$ and relations
\begin{equation}\label{eq:flat-relations}
 \left[t_H,\sum_{K\supset T}t_K\right]=0
 \qquad(H\supset T),
\end{equation}
where $T$ ranges over the codimension-two intersections of the
arrangement and the sum ranges over its hyperplanes containing $T$.
This is the level-two cyclotomic infinitesimal braid arrangement
\cite[Section~1.4]{Enriquez}.  We use the integral incidence
presentation \eqref{eq:flat-relations}, not a characteristic-zero
freeness assertion.
All intersections in this definition are taken over $\Q$.
Set $\mathfrak h_d=\mathfrak h_{\Z,d}\otimes\F$, and write
\[
 \begin{gathered}
 a_i=t_{\{z_i=0\}},\qquad
 u_{ij}^s=t_{\{z_i-sz_j=0\}}\quad(s\in\{+1,-1\}),\\
 u_{ij}=u_{ij}^{+1},\quad v_{ij}=u_{ij}^{+1}+u_{ij}^{-1}.
 \end{gathered}
\]
Thus plus and minus hyperplanes remain distinct generators even after
scalar reduction.  For $d=3$ abbreviate these Lie algebras to
$\mathfrak h_\Z$ and $\mathfrak h$.  Each target symbol $a_i,v_{ij},u_{ij}$
has enveloping degree one.

The \emph{count} of a word in either named alphabet is its number of
$u$ symbols; $a,v$ symbols have count zero.  Count is attached to an
expression in the generators, not asserted to be a grading of either
quotient algebra.  We can now state the relation to be transferred.

\begin{proposition}\label{prop:leading-pentagon}
Let $\bar\psi\in\overline L_n$, $n>2$, and suppose its expression
$\Psi(A,C,B)$ lies in $F^r$.  Let $\chi$ be its component of count $r$,
and write $n=2m+r$, with $m>0$.  In $U(\mathfrak h)$ let
$\mathcal U_{\ell,<r}$ be the span of all products of exactly $\ell$
symbols from $\{a_i,v_{ij},u_{ij}\}$ that contain fewer than $r$
symbols $u$.  Then
\begin{equation}\label{eq:leading-pentagon}
\begin{split}
 0\equiv{}&\chi(a_1,v_{12}+v_{13},u_{12}+u_{13})
       +\chi(a_2+v_{12},v_{23},u_{23})\\
 &+\chi(a_1+a_2+v_{12},v_{13}+v_{23},u_{13}+u_{23})\\
 &+\chi(a_1,v_{12},u_{12})\pmod{\mathcal U_{m+r,<r}}.
\end{split}
\end{equation}
For $r=0$ the indicated lower-count subspace is zero.
\end{proposition}

The lower-count remainder belongs to the stated linear subspace of
words of fixed target length $m+r$.  In the next section, scalar deletion
tests on words of length $m+r$ containing $r$ zeros will annihilate
this subspace.
To prove the proposition, we compare
ordering calculations in the two alphabets; no homomorphism between
the two algebras is used.  In the source, terms omitted from the leading
pentagon have higher count; in the target, the permitted remainder has
lower count.  We will show that source ordering never lowers count and
target ordering never raises it, with identical count-preserving parts.

Add a newest strand $d$.  The old symbols
are called base symbols; the new symbols are called fiber symbols.
In the source write
\[
 X=X_d,\qquad H=X^2,\qquad
 U_i=u_{id},\qquad V_i=[X,U_i]\quad(i<d),
\]
and in the target write
\[
 H=a_d,\qquad U_i^s=u_{id}^s,\qquad
 U_i=U_i^{+1},\qquad V_i=U_i^{+1}+U_i^{-1}.
\]
In this notation count includes the $U$ symbols, while $H,V$ have
count zero.  The following exact formulas will justify the two
opposite count inequalities.

\begin{lemma}\label{lem:action-tables}
All commutators of a base symbol with a fiber symbol are given by the
following formulas.  In both algebras use the abbreviations
\[
 R_i=[H,U_i],\qquad Q_i=[H+V_i,U_i],\qquad
 W_{ij}=[V_i,U_j]+[U_i,V_j].
\]
In the source the exact formulas are
\begin{align}
 [a_i,H]&=[H,V_i]+[R_i,U_i],\notag\\
 [a_i,U_i]&=Q_i,& [a_i,U_k]&=0\quad(k\ne i),\notag\\
 [a_i,V_i]&=[H,V_i]+[[V_i,U_i],U_i],\notag\\
 [a_i,V_k]&=[Q_i,U_k]\quad(k\ne i),
 \label{eq:source-a-actions}
\end{align}
\begin{align}
 [u_{ij},H]&=0,\notag\\
 [u_{ij},U_i]&=[u_{ij},U_j]=[U_i,U_j],\notag\\
 [u_{ij},V_i]&=[u_{ij},V_j]=W_{ij},\notag\\
 [u_{ij},U_k]&=[u_{ij},V_k]=0\quad(k\notin\{i,j\}),
 \label{eq:source-u-actions}
\end{align}
and
\begin{align}
 [v_{ij},H]&=[H,[U_i,U_j]],\notag\\
 [v_{ij},U_i]&=[U_i,V_j],&
 [v_{ij},U_j]&=[U_j,V_i],\notag\\
 [v_{ij},V_i]&=[V_i,V_j]+[W_{ij},U_i]+[U_i,R_j],\notag\\
 [v_{ij},V_j]&=[V_i,V_j]+[W_{ij},U_j]+[U_j,R_i],\notag\\
 [v_{ij},U_k]&=0\quad(k\notin\{i,j\}),\notag\\
 [v_{ij},V_k]&=[W_{ij},U_k]\quad(k\notin\{i,j\}).
 \label{eq:source-v-actions}
\end{align}
In the target the exact formulas are
\begin{align}
 [a_i,H]&=[H,V_i],& [a_i,U_i]&=Q_i,\notag\\
 [a_i,V_i]&=[H,V_i],&
 [a_i,U_k]&=[a_i,V_k]=0\quad(k\ne i),
 \label{eq:target-a-actions}
\end{align}
\begin{align}
 [u_{ij},H]&=0,\notag\\
 [u_{ij},U_i]&=[u_{ij},U_j]=[U_i,U_j],\notag\\
 [u_{ij},V_i]&=[u_{ij},V_j]=W_{ij}+[V_i,V_j],\notag\\
 [u_{ij},U_k]&=[u_{ij},V_k]=0\quad(k\notin\{i,j\}),
 \label{eq:target-u-actions}
\end{align}
and
\begin{align}
 [v_{ij},H]&=0,\notag\\
 [v_{ij},U_i]&=[U_i,V_j],&
 [v_{ij},U_j]&=[U_j,V_i],\notag\\
 [v_{ij},V_i]&=[v_{ij},V_j]=[V_i,V_j],\notag\\
 [v_{ij},U_k]&=[v_{ij},V_k]=0\quad(k\notin\{i,j\}).
 \label{eq:target-v-actions}
\end{align}
Every source output has at least the count of its input pair; every
target output has at most that count.  The parts preserving count
are identical in the two lists and have symbolic length two.
\end{lemma}

\begin{proof}
For the source, the braid relations give the following derivations
on its free fiber on $X,U_1,\ldots,U_{d-1}$:
\[
 \begin{aligned}
 \delta_i=\ad_{X_i}:&\quad
 X\longmapsto V_i,\quad U_i\longmapsto V_i,\quad
 U_k\longmapsto0\ (k\ne i),\\
 \epsilon_{ij}=\ad_{u_{ij}}:&\quad
 X\longmapsto0,\quad U_i,U_j\longmapsto[U_i,U_j],\quad
 U_k\longmapsto0\ (k\notin\{i,j\}).
 \end{aligned}
\]
The derivation rule and $[X,V_i]=[X^2,U_i]=R_i$ imply
\[
 \delta_iH=R_i,\qquad \delta_iV_i=Q_i,\qquad
 \delta_iV_k=[V_i,U_k]\quad(k\ne i).
\]
Since $\ad_{a_i}=\delta_i^2$, applying $\delta_i$ again gives
\eqref{eq:source-a-actions}.  For its incident $V_i$ entry, in
particular, $R_i+Q_i=[V_i,U_i]$ gives
\[
 \delta_iQ_i=[R_i+Q_i,U_i]+[H+V_i,V_i]
 =[ [V_i,U_i],U_i]+[H,V_i].
\]
Applying $\epsilon_{ij}$ to $H,V_k$ gives
\eqref{eq:source-u-actions}, since
$[X,[U_i,U_j]]=W_{ij}$.

Next $\ad_{v_{ij}}=[\delta_i,\epsilon_{ij}]$ sends $X$ to $W_{ij}$,
sends $U_i$ to $[U_i,V_j]$, sends $U_j$ to $[U_j,V_i]$, and kills the
other $U_k$.  Its action on $H$ is
$[X,W_{ij}]=[H,[U_i,U_j]]$.  Its remaining actions follow uniformly
from
\[
 [v_{ij},V_k]=[W_{ij},U_k]+[X,[v_{ij},U_k]],
\]
which proves \eqref{eq:source-v-actions}, including the nonincident
entries.

For the target, the flat where $z_i=z_d=0$ lies in exactly the four
hyperplanes with generators $a_i,H,U_i^{+1},U_i^{-1}$.  Its relations give
\[
 [a_i,H]=[H,V_i],\qquad
 [a_i,U_i^t]=[H+V_i,U_i^t]\quad(t\in\{+1,-1\}).
\]
Summing the second formula over $t$ proves the $V_i$ entry of
\eqref{eq:target-a-actions}.  The signed triple flat containing
$u_{ij}^s,U_i^t,U_j^{st}$ gives
\[
 [u_{ij}^s,U_i^t]=[U_j^{st},U_i^t],\qquad
 [u_{ij}^s,U_j^t]=[U_i^{st},U_j^t].
\]
Substitute $U_k^{-1}=U_k+V_k$.  For example,
\[
 \begin{aligned}
 [u_{ij},V_i]
 &=[U_j,U_i]+[U_j+V_j,U_i+V_i]\\
 &=W_{ij}+[V_i,V_j].
 \end{aligned}
\]
Summing also over $s$ gives
$[v_{ij},U_i]=[U_i,V_j]$ and
$[v_{ij},V_i]=[V_i,V_j]$; exchanging $i,j$ gives the other incident
entries.  Every remaining pair in the claimed zero entries meets
in a flat contained in just those two hyperplanes, so its generators
commute.  This proves all target formulas.

The count assertions can now be read term by term.  In the source,
the terms beyond the common quadratic terms add two $U$'s.  In the
target, the only extra term is $[V_i,V_j]$ in
\eqref{eq:target-u-actions}, whose count is zero instead of one.
This also proves agreement of the parts preserving count.
\end{proof}

For clarity, their common incident actions are the short table
\begin{equation}\label{eq:common-actions}
\begin{array}{c|ccc}
 &H&U_i&V_i\\ \hline
 a_i &[H,V_i]&[H+V_i,U_i]&[H,V_i]\\
 u_{ij}&0&[U_i,U_j]&[U_i,V_j]+[V_i,U_j]\\
 v_{ij}&0&[U_i,V_j]&[V_i,V_j]
\end{array}
\end{equation}
with the exchanged-index versions supplied by the full formulas
above and zero common actions at nonincident indices.

We next specify exactly how to use the tables.  At a fixed newest
strand, scan a word from left to right for its first adjacent pair
$bf$ consisting of a base symbol followed by a fiber symbol.  Replace
it by
\begin{equation}\label{eq:ordering-swap}
 bf=fb+[b,f],
\end{equation}
using the appropriate exact formula of
Lemma~\ref{lem:action-tables} and expanding brackets into associative
words.  Fix the order of those expansions once and for all.
On each resulting branch the lexicographic pair
\[
 \bigl(\text{number of base symbols},\,
        \text{number of base-before-fiber inversions}\bigr)
\]
strictly decreases.  The swapped branch lowers the second entry;
every commutator branch removes a base symbol and inserts only fiber
symbols, hence lowers the first entry.  This proves termination even
for a source correction that increases symbolic length.  When all
fiber letters precede the base suffix, repeat the procedure on that
suffix with one fewer strand.  The result is an iterated ordered
expression with the newest fiber first.

There are two further facts about this algorithm.  First, in the
source its ordered products are linearly independent.  Indeed,
Lemma~\ref{lem:free-fiber} identifies the ordinary enveloping algebra,
as a vector space, with the tensor product of the free associative
fiber algebras in this order.  In each fiber the subalgebra generated
by $H,V_i,U_i$ is free associative: with $X$ larger than all $U_i$,
these elements have the respective leading words
\[
 XX,\qquad XU_i,\qquad U_i.
\]
They again form a prefix code.  Thus distinct words in each restricted
fiber alphabet are independent, and so are their ordered products
over all fibers.  At the final one-strand stage this says simply
that the powers of $X_1^2$ are independent.

Second, source ordering never lowers count, while target ordering
never raises it.  This remains true through every recursive stage,
since the unchanged surrounding word contributes the same count
and the action formulas have the required inequality term by term.
A source branch that has acquired count greater than $r$ can never
contribute at count $r$; a target branch of count less than $r$ can
never return to count $r$.  Consequently the branches that preserve
count throughout are described in both algebras by exactly the same
algorithm using \eqref{eq:common-actions}.

\begin{proof}[Proof of Proposition~\ref{prop:leading-pentagon}]
Relabel the four braid strands $0,1,2,3$.  Since signs disappear over
$\F$, the reduced pentagon says that the sum of
\[
 \begin{gathered}
 \bar\psi(X_1,u_{12}+u_{13}),\qquad
 \bar\psi(X_2+u_{12},u_{23}),\qquad
 \bar\psi(u_{12},u_{23}),\\
 \bar\psi(X_1+X_2,u_{13}+u_{23}),\qquad
 \bar\psi(X_1,u_{12})
 \end{gathered}
\]
is zero in the source.  Express these five terms using $\Psi$.
Their exact triples of arguments $(A,C,B)$ are, in the same order,
\begin{equation}\label{eq:five-substitutions}
\begin{aligned}
 &(a_1,\ v_{12}+v_{13},\ u_{12}+u_{13}),\\
 &(a_2+v_{12}+u_{12}^2,\ v_{23}+[u_{12},u_{23}],\ u_{23}),\\
 &(u_{12}^2,\ [u_{12},u_{23}],\ u_{23}),\\
 &(a_1+a_2+v_{12},\ v_{13}+v_{23},\ u_{13}+u_{23}),\\
 &(a_1,\ v_{12},\ u_{12}).
\end{aligned}
\end{equation}
In particular,
\[
 (X_2+u_{12})^2=a_2+v_{12}+u_{12}^2,
 \qquad (X_1+X_2)^2=a_1+a_2+v_{12}.
\]
Only the former square has a $u_{12}^2$ summand.

Expand the pentagon as a polynomial in the restricted source symbols
before ordering.  Its part of count exactly $r$ is precisely the
four-term polynomial $P$ on the right of
\eqref{eq:leading-pentagon}.  To verify this, the higher components
of $\Psi$ already have more than $r$ occurrences of $B$, each of
which becomes a single $u$.  Either extra summand $u_{12}^2$ or
$[u_{12},u_{23}]$ in the second triple increases count by two.
The third triple, being entirely in the $u$ alphabet, gives count
$n>r$.  All remaining replacements preserve the original $B$-count.
Thus the full source pentagon has the form $P+E=0$, where every word
of $E$ has count greater than $r$.

The distinction between original weight and symbolic length is
essential here.  Assign source weight two to each $a,v,H,V$ and
weight one to each $u,U$.  A word of length $\ell$ and count $q$ has
original weight
\begin{equation}\label{eq:weight-length-count}
 n=2\ell-q.
\end{equation}
All source action formulas preserve this weight.  Thus a correction
that raises count by two raises length by one, and every ordered
word of weight $n$ and count $r$ has length
$(n+r)/2=m+r$.  In particular $P$ has that length, while the
higher-count expressions are allowed to have greater length.

Apply the deterministic source ordering to $P+E=0$.  Since count
never decreases, $E$ contributes nothing at count $r$.  Independence
of the ordered restricted source products, proved above from the
free fibers, implies that every ordered coefficient of count $r$
in $P$ is zero separately.  This conclusion uses independence
across all the restricted words, so expressions of other lengths
or counts cannot cancel a coefficient under consideration.

Now interpret only the polynomial $P$, of length $m+r$, in the
target.  Apply the same deterministic ordering, discarding a branch
as soon as its count falls below $r$.  Such a branch can never return
to count $r$.  The retained branches use the identical quadratic
formulas that retained count $r$ in the source.  Their ordered
coefficient vector is therefore the zero vector just obtained.
Every target replacement preserves length, so every discarded
branch has length $m+r$ and count less than $r$.  Their sum belongs
to $\mathcal U_{m+r,<r}$, proving \eqref{eq:leading-pentagon}.

Only source independence enters this argument.  In the target the
ordering procedure is a sequence of valid equalities yielding a
spanning expression; no independence or free-fiber assertion for
$\mathfrak h$ is required.
\end{proof}
